\section{Model: Support Vector Machine (SVM)}

\subsection{Introduction}
This report analyzes the performance of the Support Vector Machine (SVM) model trained using various embedding methods. The model was implemented using the \texttt{SVC} class from \texttt{scikit-learn} with different kernel functions and hyperparameter configurations. The primary objective was to achieve high classification accuracy while maintaining robust generalization across different embedding techniques.

\subsection{Training Configuration}
The SVM model was trained with the following hyperparameter search space:
\begin{itemize}
    \item \textbf{Kernel}: \texttt{linear}, \texttt{rbf}, \texttt{poly}, \texttt{sigmoid}.
    \item \textbf{Regularization Parameter} (\texttt{C}): 0.1, 1.0, 10.0.
    \item \textbf{Gamma}: \texttt{scale}, \texttt{auto}.
\end{itemize}

The best hyperparameters selected based on model evaluation were:
\begin{itemize}
    \item \textbf{Kernel}: \texttt{linear}
    \item \textbf{C}: 0.1 (for Count, Word2Vec, GloVe), 1.0 (for TF-IDF)
    \item \textbf{Gamma}: \texttt{scale}
\end{itemize}

\subsection{Training and Evaluation Results}
The model was trained and evaluated using K-Fold Cross-Validation across different feature extraction methods: Count Vectorizer, TF-IDF, Word2Vec, and GloVe. The best model was selected based on Accuracy, with secondary considerations for F1-score and ROC AUC.

\textbf{Training Performance Metrics:}

\begin{table}[H]
    \centering
    \caption{Training Performance Metrics for SVM}
    \label{tab:svm-training-metrics}
    \begin{tabular}{|l|c|c|c|c|c|}
        \hline
        \textbf{Method} & \textbf{Accuracy} & \textbf{ROC AUC} & \textbf{F1} & \textbf{Precision} & \textbf{Recall} \\ 
        \hline
        Count Vectorizer & 0.71 & 0.70 & 0.73 & 0.68 & 0.79 \\ 
        \hline
        TF-IDF & 0.70 & 0.70 & 0.71 & 0.69 & 0.74 \\ 
        \hline
        Word2Vec & 0.71 & 0.71 & 0.72 & 0.71 & 0.74 \\ 
        \hline
        GloVe & 0.71 & 0.71 & 0.71 & 0.71 & 0.72 \\ 
        \hline
    \end{tabular}
\end{table}

\textbf{Testing Performance Metrics:}

\begin{table}[H]
    \centering
    \caption{Testing Performance Metrics for SVM}
    \label{tab:svm-testing-metrics}
    \begin{tabular}{|l|c|c|c|c|c|}
        \hline
        \textbf{Method} & \textbf{Accuracy} & \textbf{ROC AUC} & \textbf{F1} & \textbf{Precision} & \textbf{Recall} \\ 
        \hline
        Count Vectorizer & 0.7060 & 0.7770 & 0.7325 & 0.6857 & 0.7861 \\ 
        \hline
        TF-IDF & 0.6880 & 0.7663 & 0.7040 & 0.6845 & 0.7246 \\ 
        \hline
        Word2Vec & 0.7160 & 0.7846 & 0.7290 & 0.7127 & 0.7461 \\ 
        \hline
        GloVe & 0.6945 & 0.7693 & 0.7030 & 0.6999 & 0.7061 \\ 
        \hline
    \end{tabular}
\end{table}

\textbf{Best Model Selection Criteria:}

\begin{itemize}
    \item The best model is chosen based on testing performance rather than training performance.
    \item The selection priority follows: Accuracy > F1 Score > ROC AUC.
    \item Based on this criterion, the best model is:
\end{itemize}

\begin{verbatim}
{
    "method": "word2vec",
    "model": "svm",
    "hyperparameters": { "C": 0.1, "gamma": "scale", "kernel": "linear" },
    "performance": {
        "accuracy": 0.7160,
        "precision": 0.7127,
        "recall": 0.7461,
        "f1": 0.7290,
        "roc_auc": 0.7846
    }
}
\end{verbatim}

\textbf{Conclusion:} The SVM model trained with Word2Vec embeddings achieved the highest accuracy (0.7160) and the best overall balance across F1-score and ROC AUC, making it the optimal choice for sentiment classification in our experiment when it comes to SVM.

\subsection{Performance Analysis}
\begin{itemize}
    \item \textbf{Accuracy Analysis}: The model trained on Word2Vec achieved the highest accuracy (71.60\%), outperforming Count Vectorizer, TF-IDF, and GloVe embeddings.
    \item \textbf{Kernel Effectiveness}: The linear kernel consistently performed best across all embedding methods.
    \item \textbf{ROC AUC}: The model exhibited strong class differentiation with ROC AUC scores ranging from 76.63\% to 78.46\%.
    \item \textbf{Precision and Recall}: The Word2Vec model maintained a good balance with precision of 71.27\% and recall of 74.61\%.
    \item \textbf{Embedding Comparison}: Word2Vec performed best, while GloVe showed the lowest performance among embedding methods.
\end{itemize}

\subsection{Visualization of Training Results}
The following figures illustrate the model's performance across different embedding techniques:

\begin{figure}[H]
    \centering
    \begin{subfigure}[b]{0.48\textwidth}
        \includegraphics[width=\textwidth]{img/report_info/img/9_3_SVM/best_svm_count_loss.png}
        \caption{Loss Curve - Count Vectorizer}
        \label{fig:svm-count-loss}
    \end{subfigure}
    \begin{subfigure}[b]{0.48\textwidth}
        \includegraphics[width=\textwidth]{img/report_info/img/9_3_SVM/best_svm_tfidf_loss.png}
        \caption{Loss Curve - TF-IDF}
        \label{fig:svm-tfidf-loss}
    \end{subfigure}
    
    \begin{subfigure}[b]{0.48\textwidth}
        \includegraphics[width=\textwidth]{img/report_info/img/9_3_SVM/best_svm_word2vec_loss.png}
        \caption{Loss Curve - Word2Vec}
        \label{fig:svm-word2vec-loss}
    \end{subfigure}
    \begin{subfigure}[b]{0.48\textwidth}
        \includegraphics[width=\textwidth]{img/report_info/img/9_3_SVM/best_svm_glove_loss.png}
        \caption{Loss Curve - GloVe}
        \label{fig:svm-glove-loss}
    \end{subfigure}
    
    \caption{Comparison of Loss Curves for SVM across Different Feature Extraction Methods}
    \label{fig:svm-loss-group}
\end{figure}

\begin{figure}[H]
    \centering
    \begin{subfigure}[b]{0.48\textwidth}
        \includegraphics[width=\textwidth]{img/report_info/img/9_3_SVM/best_svm_count.png}
        \caption{Performance - Count Vectorizer}
        \label{fig:svm-count}
    \end{subfigure}
    \begin{subfigure}[b]{0.48\textwidth}
        \includegraphics[width=\textwidth]{img/report_info/img/9_3_SVM/best_svm_tfidf.png}
        \caption{Performance - TF-IDF}
        \label{fig:svm-tfidf}
    \end{subfigure}
    
    \begin{subfigure}[b]{0.48\textwidth}
        \includegraphics[width=\textwidth]{img/report_info/img/9_3_SVM/best_svm_word2vec.png}
        \caption{Performance - Word2Vec}
        \label{fig:svm-word2vec}
    \end{subfigure}
    \begin{subfigure}[b]{0.48\textwidth}
        \includegraphics[width=\textwidth]{img/report_info/img/9_3_SVM/best_svm_glove.png}
        \caption{Performance - GloVe}
        \label{fig:svm-glove}
    \end{subfigure}
    
    \caption{Comparison of Training Performance Metrics for SVM across Different Feature Extraction Methods}
    \label{fig:svm-performance-group}
\end{figure}

\textbf{Image Description:}

\begin{itemize}
    \item \textbf{Training and Validation Loss Analysis:}
    \begin{itemize}
        \item Word2Vec shows the most stable validation loss across folds.
        \item Count Vectorizer exhibits higher variance in validation performance.
        \item All methods show convergence in training loss.
    \end{itemize}
    
    \item \textbf{Validation Performance Metrics:}
    \begin{itemize}
        \item Word2Vec achieves the highest and most consistent accuracy and ROC AUC.
        \item Count Vectorizer shows higher recall but lower precision.
        \item TF-IDF demonstrates moderate but stable performance.
        \item GloVe has the lowest performance consistency across metrics.
    \end{itemize}
\end{itemize}

\subsection{Conclusion}

The SVM model performed best with Word2Vec embeddings, achieving an accuracy of 71.60\%. The linear kernel proved most effective across all embedding methods. While the model showed good generalization (ROC AUC of 78.46\%), its overall performance was slightly lower than Logistic Regression. Future improvements could include: Experimenting with non-linear kernels and custom embedding dimensions...

Overall, the SVM model with Word2Vec features provides a robust alternative for sentiment classification, particularly when interpretability of linear decision boundaries is desired.